\section{Occupation bounds for the intrinsic constant}
\label{S:intrinsic-occupation-ceiling}

Retain all original models and amplitudes, put $a=\alpha_0$,
$q=4-\alpha>2$, $k=3-\alpha=q-1>1$ and $m=q/k$.
The temporal kernel is $K(s,t)=|s-t|^{-a}$ on $I=(0,1)$;
its diagonal value may be set to zero.
Let $\mathcal O$ be the full strongly $H^1$-measurable unit-profile
class defined before Proposition~\ref{P:occupation-physical-class}.
Here we use the \emph{raw} kernel, defining
\begin{equation}\label{E:raw-occupation-definitions}
  \begin{aligned}
    A(g)&=\int_{I^2}K(s,t)B_\Gamma(g_s^2,g_t^2)\ud s\ud t,\\
    C_{\rm occ}&=\sup_{g\in\mathcal O}
    \frac{A(g)}{T(g)^{\alpha/2}},&
    \mathcal E_\zeta^{\mathcal O}
    &=\sup_{g\in\mathcal O}\{A(g)-\zeta T(g)\},\qquad \zeta>0.
  \end{aligned}
\end{equation}
The scalar penalty $\zeta$ is not a spatial mask.
These objects differ from the earlier damped-kernel energy.
Assume $\Gamma\ne0$ until the direct zero case below.

\begin{lemma}[Full-profile squares in the original raw space]
  \label{L:occupation-raw-square}
  There is a continuous map $Q:H^1(\R^d;\R)\to H_{\mathrm{sp}}$ with
  \[
    \langle Q(f),Q(h)\rangle_{H_{\mathrm{sp}}}
    =B_\Gamma(f^2,h^2),\qquad
    \|Q(f)\|^2=A_{\mathrm{sp}}(f).
  \]
  For $g\in\mathcal O$, the element
  $\rho_g=J_{H_{\mathrm{sp}}}(s\mapsto Q(g_s))$ belongs to the ORIGINAL
  $H_I$ and satisfies
  \begin{equation}\label{E:occupation-raw-identities}
    \|\rho_g\|_{H_I}^2=A(g),\qquad
    \langle\Psi,\rho_g\rangle_{H_I}
    =\int_I\int_{\R^d}h_{\Psi,s}(x)g_s(x)^2\ud x\ud s
  \end{equation}
  for every $\Psi\in\mathcal C_{\mathrm{pos}}$. In particular,
  \[
    \langle\Psi,\rho_g\rangle_{H_I}\le\sqrt{v_\Psi A(g)}.
  \]
\end{lemma}

\begin{proof}
  Approximate $f$ in $H^1$ by real compact smooth $f_n$.
  Proposition~\ref{P:spatial-four-linear} gives
  \[
    \|f_n^2-f_l^2\|_{H_{\mathrm{sp}}}^2
    =C_4(f_n-f_l,f_n+f_l,f_n-f_l,f_n+f_l)
    \le R_{\mathrm{sp}}N(f_n-f_l)^2N(f_n+f_l)^2\to0.
  \]
  Here $N$ is continuous on $H^1$ and $N(f)\le\|f\|_{H^1}$.
  Completeness defines $Q(f)$; interlacing two approximating sequences
  proves independence. Passing the physical four-linear continuity
  and Hilbert pairings to limits proves the first identities.
  The same difference bound, first for smooth approximants and then
  by passage to the limit, proves continuity of $Q$.

  Thus $F_g(s)=Q(g_s)$ is strongly measurable.
  Writing $b_s=\|\nabla g_s\|_2^2$, the unit norm and
  $0<\alpha/2<1$ give
  \[
    \int_I\|F_g(s)\|^2\ud s
    \le R_{\mathrm{sp}}\int_I b_s^{\alpha/2}\ud s
    \le R_{\mathrm{sp}}T(g)^{\alpha/2}<\infty .
  \]
  Lemma~\ref{L:positive-raw-embedding} now constructs $\rho_g$ and
  identifies its squared norm with the physical $A(g)$.
  Joint measurability and representative independence follow from
  the continuous physical pairing on $(H^1)^2$.
  Absolute time integrability follows from the same Schur bound.

  Choose one compact smooth $\chi$ equal to one on all spatial
  supports of $\Psi_t$. For real $f\in H^1$, smooth approximation gives
  \[
    \langle\Psi_t,Q(f)\rangle_{H_{\mathrm{sp}}}
    =\lim_n B_\Gamma(\Psi_t,f_n^2)
    =C_4(\chi,\Psi_t,f,f)=B_\Gamma(\Psi_t,f^2).
  \]
  Proposition~\ref{P:spatial-four-linear} justifies these absolute
  integrals and limits.
  Polarizing~\eqref{E:positive-bochner-inclusion} gives
  \[
    \langle\Psi,\rho_g\rangle
    =\int_{I^2}K(s,t)\langle\Psi_t,Q(g_s)\rangle_{H_{\mathrm{sp}}}
    \ud t\ud s .
  \]
  The absolute integrand is bounded by the product of two
  $H_{\mathrm{sp}}$ norms against $K$, integrable by Schur and
  $L^2$ Cauchy--Schwarz.
  Evenness of $\Gamma$ and Tonelli for the nonnegative physical
  integrands yield~\eqref{E:occupation-raw-identities}.
  Hilbert Cauchy--Schwarz with $\|\Psi\|^2=v_\Psi$ proves the bound.
  No rough occupation is inserted into a stochastic flow.
  This proves Lemma~\ref{L:occupation-raw-square}.
\end{proof}

\begin{proposition}[Raw occupation energy and dilation]
  \label{P:raw-occupation-energy}
  With $K_0=2^a/(1-a)$,
  \[
    0<C_{\rm occ}\le K_0R_{\mathrm{sp}}<\infty .
  \]
  Moreover
  \begin{align}
    \mathcal E_\zeta^{\mathcal O}
    &=D_\zeta C_{\rm occ}^{2/(2-\alpha)},&
    D_\zeta&=\frac{2-\alpha}{2}
    \left(\frac{\alpha}{2\zeta}\right)^{\alpha/(2-\alpha)},
    \label{E:raw-occupation-energy}\\
    C_{\rm occ}
    &=\left(\frac{2\mathcal E_\zeta^{\mathcal O}}{2-\alpha}\right)^{(2-\alpha)/2}
    \left(\frac{2\zeta}{\alpha}\right)^{\alpha/2}.
    \label{E:raw-occupation-ratio}
  \end{align}
  In particular $\mathcal E_\zeta^{\mathcal O}>0$ is finite and
  equals $\zeta^{-\alpha/(2-\alpha)}\mathcal E_1^{\mathcal O}$.
\end{proposition}

\begin{proof}
  A unit $H^1$ profile has positive gradient energy:
  a zero weak gradient makes it constant, and an $L^2$ constant
  on $\R^d$ is zero. Hence $0<T(g)<\infty$.
  Spatial covariance Cauchy--Schwarz gives
  $B_\Gamma(g_s^2,g_t^2)\le
  R_{\mathrm{sp}}b_s^{\alpha/4}b_t^{\alpha/4}$.
  Schur with row bound $K_0$, then concavity, proves
  \begin{equation}\label{E:raw-occupation-bound}
    0\le A(g)\le K_0R_{\mathrm{sp}}T(g)^{\alpha/2}.
  \end{equation}
  A fixed compact smooth unit profile, constant in time, belongs to
  $\mathcal O$ and has positive $A$ by
  Proposition~\ref{P:spatial-four-linear}. This proves positivity
  and finiteness of $C_{\rm occ}$.

  The map $(U_\ell g)_s(x)=\ell^{d/2}g_s(\ell x)$ is an $H^1$
  isomorphism preserving unit mass and strong measurability.
  It has inverse $U_{1/\ell}$ and is a bijection on $\mathcal O$.
  The physical Radon change of variables gives
  \[
    A(U_\ell g)=\ell^\alpha A(g),\qquad
    T(U_\ell g)=\ell^2T(g).
  \]
  Indeed the two densities contribute $\ell^{2d}$, $\ud x$ contributes
  $\ell^{-d}$, and the pushforward of $\Gamma$ contributes
  $\ell^{\alpha-d}$. Approximation extends the identity to full $H^1$.
  Taking the union over $g,\ell$ leaves the class unchanged.
  For fixed $A,T>0$, the derivative of
  $A\ell^\alpha-\zeta T\ell^2$ has the sign of
  $\alpha A-2\zeta T\ell^{2-\alpha}$.
  Its global maximum occurs at
  $\ell=(\alpha A/(2\zeta T))^{1/(2-\alpha)}$ and equals
  $D_\zeta[A/T^{\alpha/2}]^{2/(2-\alpha)}$.
  For $A=0$ the supremum is zero as $\ell\downarrow0$.
  Bounding every ratio by $C_{\rm occ}$ gives one inequality in
  \eqref{E:raw-occupation-energy}.
  Fixed ratios approaching its positive supremum give the other
  by continuity of the positive power.
  Algebra gives~\eqref{E:raw-occupation-ratio} and the penalty scaling.
  No optimizer or minimax interchange is assumed.
  This proves Proposition~\ref{P:raw-occupation-energy}.
\end{proof}

\begin{lemma}[Countable Rayleigh near-occupations]
  \label{L:occupation-rayleigh-selection}
  For fixed $\Psi\in\mathcal C_{\mathrm{pos}}$ and every $\varepsilon>0$
  there is $g_\varepsilon\in\mathcal O$ such that
  \begin{equation}\label{E:occupation-rayleigh-selection}
    I_\Psi^2+T(g_\varepsilon)
    \le\sqrt{v_\Psi A(g_\varepsilon)}+\varepsilon .
  \end{equation}
\end{lemma}

\begin{proof}
  The real $H^1$ space is separable, for example by its isometric
  embedding $f\mapsto(f,\nabla f)$ into a finite product of $L^2$ spaces.
  Compact smooth density therefore supplies a countable real family.
  Discard zero and normalize each member in $L^2$ to obtain $f_j$
  dense in the $L^2$-unit $H^1$ class. Normalization preserves convergence
  because the approximating $L^2$ norms tend to one.
  Put $b_j=\|\nabla f_j\|_2^2$ and
  $R_j(s)=\int h_{\Psi,s}f_j^2-b_j$.
  Bounded-potential form continuity and
  Proposition~\ref{P:variational-unmasked-rate} give
  $\lambda_{\Psi,s}^2=\sup_jR_j(s)$.
  The canonical functions are measurable, indeed continuous.
  Let $U_j=\{s:R_j(s)>\lambda_{\Psi,s}^2-\varepsilon\}$,
  $A_1=U_1$ and $A_j=U_j\setminus\bigcup_{i<j}U_i$.
  These first-index cells are Borel and cover $I$.
  Put $g_\varepsilon(s)=f_j$ on $A_j$.
  Finite-valued truncations using $f_1$ on the remaining cells
  converge pointwise in $H^1$, proving strong measurability.
  A null-set representative may be filled by $f_1$.

  With $H_\Psi=\sup_{s,x}h_{\Psi,s}(x)<\infty$, the selection implies
  \[
    b(g_\varepsilon(s))
    <\int h_{\Psi,s}g_\varepsilon(s)^2-\lambda_{\Psi,s}^2+\varepsilon
    \le H_\Psi+\varepsilon .
  \]
  Thus it belongs to the actual $\mathcal O$.
  Integrating the near-Rayleigh inequality, using
  Lemma~\ref{L:occupation-raw-square}, gives
  \[
    \int_I\lambda_{\Psi,s}^2\ud s+T(g_\varepsilon)
    \le\langle\Psi,\rho_{g_\varepsilon}\rangle+\varepsilon
    \le\sqrt{v_\Psi A(g_\varepsilon)}+\varepsilon .
  \]
  Time Cauchy--Schwarz on the unit interval gives
  $I_\Psi^2\le\int_I\lambda_{\Psi,s}^2\ud s$.
  This proves~\eqref{E:occupation-rayleigh-selection}.
  No eigenvector, continuum selector or limiting occupation is required.
  This proves Lemma~\ref{L:occupation-rayleigh-selection}.
\end{proof}

For $0\le b<q/2$ define $F_b(0)=0$ and
\begin{equation}\label{E:occupation-scalar-map}
  F_b(z)=k\left[
  \frac{(q-2b)^{(q-2b)/2}}{8(q-b)^{q-b}}
  \left(\frac{2z}{2-\alpha}\right)^{(2-\alpha)/2}
  \right]^{1/k},\qquad z>0.
\end{equation}
At $b=0$ this is exactly the map in~\eqref{E:spatial-static-energy-score}.
The display parameter $b$ does not change the actual temporal exponent $a$.

\begin{theorem}[The computed occupation ceiling]\label{T:intrinsic-occupation-ceiling}
  For every positive score coefficient $\mathfrak c$,
  \begin{align}
    S_{\mathfrak c}^{\mathrm{pos}}
    &\le U_{\mathfrak c}
    :=\mathfrak c^m k
    \left[\frac{\alpha^{\alpha/2}C_{\rm occ}}{8q^{q/2}}\right]^{1/k}
    \label{E:occupation-score-ceiling}\\
    &=\mathfrak c^m(2\zeta)^{\alpha/(2k)}
    F_0(\mathcal E_\zeta^{\mathcal O}),\qquad \zeta>0.\notag
  \end{align}
  In particular the actual original moment constant satisfies
  $C_2(\Gamma,a)\le U_2$ for every original model.
  This ceiling is not asserted to be attained or equal to $C_2$.
\end{theorem}

\begin{proof}
  Fix $\Psi$ and write $v=v_\Psi$, $C=C_{\rm occ}$ and $\beta=\alpha/4$.
  Lemma~\ref{L:occupation-rayleigh-selection} gives
  $I_\Psi^2\le\sqrt{vC}T^\beta-T+\varepsilon$ for some $T>0$.
  Bound it by the supremum over all $T\ge0$ before sending
  $\varepsilon\downarrow0$.
  If $v=0$, that supremum is zero.
  If $v>0$, the derivative
  $\beta\sqrt{vC}T^{\beta-1}-1$ strictly decreases from infinity to $-1$.
  Its unique positive zero is
  $T_*=(\beta\sqrt{vC})^{1/(1-\beta)}$.
  The global maximum equals
  $(q/4)(\alpha/4)^{\alpha/q}(vC)^{2/q}$.
  In both cases,
  \[
    I_\Psi\le Dv^{1/q},\qquad
    D=\sqrt{q/4}(\alpha/4)^{\alpha/(2q)}C^{1/q}.
  \]
  The scalar function $\mathfrak c Dv^{1/q}-v/2$ vanishes at zero,
  tends to minus infinity, and has strictly decreasing derivative.
  Its unique positive maximum is at
  $v_*=(2\mathfrak c D/q)^{q/k}$, with value $kv_*/2$.
  Direct positive-power algebra gives
  \[
    (2D/q)^q=q^{-q/2}(\alpha/4)^{\alpha/2}C .
  \]
  Since $k+\alpha=3$, the maximum is precisely $U_{\mathfrak c}$
  in~\eqref{E:occupation-score-ceiling}, with denominator eight.
  Substitution of~\eqref{E:raw-occupation-ratio} gives the second form,
  retaining the factor $(2\zeta)^{\alpha/(2k)}$.
  Take the supremum over the actual smooth positive-control class.
  Theorem~\ref{T:intrinsic-second-moment} identifies
  $C_2=S_2^{\mathrm{pos}}$, proving the moment ceiling.
  The penalty factor cancels its energy scaling, so the ceiling
  is independent of $\zeta$.
  This proves Theorem~\ref{T:intrinsic-occupation-ceiling}.
\end{proof}

\begin{corollary}[A common spatial-energy bracket]
  \label{C:intrinsic-spatial-bracket}
  With the same original covariance and $M_a$ of
  \eqref{E:temporal-equilibrium-mass},
  \begin{equation}\label{E:intrinsic-spatial-bracket}
    2^m M_a^{-1/k}F_0(\mathcal E_{\mathrm{sp}})
    \le C_2(\Gamma,a)
    \le2^m K_0^{1/k}F_0(\mathcal E_{\mathrm{sp}}).
  \end{equation}
  For $\Gamma\ne0$ there is also the own-energy comparison
  \[
    C_2\le2^mF_0(\mathcal E_{1/2}^{\mathcal O})
    <2^mF_a(\mathcal E_{1/2}^{\mathcal O}).
  \]
\end{corollary}

\begin{proof}
  Insert $C_{\rm occ}\le K_0R_{\mathrm{sp}}$ into
  \eqref{E:occupation-score-ceiling} and use
  \eqref{E:spatial-energy-map}.
  This gives the upper bracket; the lower is
  \eqref{E:spatial-static-energy-score} multiplied by $2^m$,
  using Theorem~\ref{T:intrinsic-second-moment}.
  Both concern one common spatial energy, not an assumed first moment.

  For completeness put
  \[
    \Delta=\tfrac12[q\log q+(q-2a)\log(q-2a)]
    -(q-a)\log(q-a)>0.
  \]
  Here $a>0$ and $q-2a>2-a>1$.
  The existing scalar strict-log-gap deduction applies at exactly
  these guards, without $q\ge3$.
  Equivalently the strict midpoint convexity of $x\log x$ on
  positive arguments, whose second derivative is $1/x$, gives the sign.
  Subtracting logarithms in~\eqref{E:occupation-scalar-map} gives
  \[
    \log(F_a(z)/F_0(z))=\Delta/k,\qquad
    \frac{\mathfrak c^mF_a(\mathcal E_\zeta^{\mathcal O})}{U_{\mathfrak c}}
    =\exp\!\left(\frac{\Delta-(\alpha/2)\log(2\zeta)}{k}\right).
  \]
  At $\zeta=1/2$ this and the ceiling prove the strict comparison.
  At $\zeta=1$ the sign is not universal.
  For $\alpha=5/4,a=1/2,d=2$,
  \[
    8[\Delta-(\alpha/2)\log2]
    =\log\frac{11^{11}7^7}{2^5 9^{18}}<0,
  \]
  since $11^{11}7^7=234966429149994773$ and
  $2^5 9^{18}=4803028329503971872$.
  This contrasts computed bounds; it refutes no actual moment inequality.
  This proves Corollary~\ref{C:intrinsic-spatial-bracket}.
\end{proof}

For $\Gamma=0$, all raw means and scores vanish and $C_2=0$.
Here $Q=0$, $A=0$ and $C_{\rm occ}=0$.
The supremum of $-\zeta T(g)$ is zero: dilate a fixed unit compact
profile by $\ell^{d/2}f(\ell x)$ and let $\ell\downarrow0$.
The same argument gives $\mathcal E_{\mathrm{sp}}=0$ and the earlier
damped energy zero, without using a zero vector as a unit profile.
With $F_b(0)=0$, every non-strict bound above is the zero identity.
These own energies are not identified with Chen's printed energies.
